\section{Biểu thức tọa độ của các phép toán véc-tơ}
\subsection{Đề 1}
\Opensolutionfile{ans}[ans/ans-2-B8]
\caulc
\begin{ex}%[Nguyễn Khánh Trọng]%[2H2N2-3]
	Trong không gian $Oxyz$, cho hai véc-tơ là $\vec{u}=(2;1;-1)$ và $\vec{v}=(1;3;1)$. Toạ độ của véc-tơ $\vec{u}+2\vec{v}$ tương ứng là
	\choice
	{$(3;4;0)$}
	{$(1;-2;-2)$}
	{\True $(4;7;1)$}
	{$(5;5;-1)$}
	\loigiai{Gọi toạ độ của véc-tơ $(\vec{u}+2\vec{v})$ là $(x;y;z)$. Khi đó
		$\heva{&x = 2 + 2\cdot1\\
			   &y = 1 + 2\cdot3\\
			   &z =  - 1 + 2\cdot1}
			   \Leftrightarrow	\heva{&x = 4\\&y = 7\\&z = 1.}$\\
	Vậy toạ độ của véc-tơ $\vec{u}+2\vec{v}$ là $(4;7;1)$.
	}
\end{ex}

\begin{ex}%[Nguyễn Khánh Trọng]%[2H2N2-2]
	Trong không gian $Oxyz$, cho điểm $A$ thỏa mãn $\overrightarrow{OA}=2\overrightarrow{i}+\overrightarrow{j}$, với $\overrightarrow{i}$, $\overrightarrow{j}$  là hai véc-tơ đơn vị trên hai trục tọa độ $Ox$, $Oy$. Tọa độ điểm $A$ là 
	\choice
	{\True $A(2;1;0)$}
	{$A(0;2;1)$}
	{$A(0;1;1)$}
	{$A(1;1;1)$}
	\loigiai{
		Ta có $\overrightarrow{OA}=2\overrightarrow{i}+\overrightarrow{j}\Rightarrow A(2;1;0)$.}
\end{ex}

\begin{ex}%[Nguyễn Khánh Trọng]%[2H2N2-4] 
	Trong không gian $Oxyz$, cho điểm $I(2;3;-4)$. Khoảng cách từ $I$ đến trục $Ox$ bằng
	\choice
	{$2$}
	{\True $5$}
	{$\sqrt{13}$}
	{$2\sqrt{5}$}
	\loigiai{Ta có $\mathrm{d}(I,Ox)=\sqrt{3^2+(-4)^2}=5$.}
\end{ex}

\begin{ex}%[Nguyễn Khánh Trọng]%[2H2N2-4] 
	Trong không gian $Oxyz$, cho hai véc-tơ $\overrightarrow{u}=(1;0;-1)$, $\overrightarrow{v}=(2;1;-2)$. Tích vô hướng $\overrightarrow{u}\cdot \overrightarrow{v}$ bằng
	\choice
	{$0$}
	{$1$}
	{\True $4$}
	{$2$}
	\loigiai{Ta có $\overrightarrow{u}\cdot \overrightarrow{v}=1\cdot 2+0\cdot 1+(-1)(-2)=4$.
	}
\end{ex}

\begin{ex}%[Nguyễn Khánh Trọng]%[2H2H2-2]
	Trong không gian $Oxyz$, cho điểm $A(2;-1;0)$ và điểm $B(3;1;1)$. Tọa độ điểm đối xứng với $A$ qua $B$ là
	\choice
	{$(1;-2;-4)$}
	{$(0; 3;-1)$}
	{\True $(4; 3; 2)$}
	{$(0;-1; 3)$}
	\loigiai{Gọi $A'$ là điểm đối xứng với $A$ qua $B$.\\
		Khi đó $B$ là trung điểm của $AA'$, ta có $\heva{& x_{A'}=2x_{B}-x_{A}=4 \\ & y_{A'}=2y_{B}-y_{A}=3 \\ & z_{A'}=2z_{B}-z_{A}=2.}$\\
		Vậy $A'(4;3;2)$.
	}
\end{ex}

\begin{ex}%[Nguyễn Khánh Trọng]%[2H2H2-2]
	Trong không gian $Oxyz$, cho tam giác $ABC$ với $A(1;2;1)$, $B(-3;0;3)$, $C(2;4;-1)$. Tìm tọa độ điểm $D$ sao cho tứ giác $ABCD$ là hình bình hành.
	\choice
	{$D(6;-6;3)$}
	{$D(6;6;3)$}
	{$D(6;-6;-3)$}
	{\True $D(6;6;-3)$}
	\loigiai{
		Gọi $D(x;y;z)$. Ta có $\overrightarrow{AB}=(-4;-2;2)$ và $\overrightarrow{DC}=(2-x;4-y;-1-z)$.\\
		Tứ giác $ABCD$ là hình bình hành $\Leftrightarrow \overrightarrow{AB}=\overrightarrow{DC}\Leftrightarrow \heva{&-4=2-x\\&-2=4-y\\&2=-1-z}\Leftrightarrow \heva{&x=6\\&y=6\\&z=-3.}$\\
		Vậy $D(6;6;-3)$.
	}
\end{ex}

\begin{ex}%[Nguyễn Khánh Trọng]%[2H2H2-3]
	Cho tam giác $ABC$ biết $A(2;-1;3)$ và trọng tâm của tam giác có toạ độ là $G(2;1;0)$. Khi đó $\overrightarrow{AB}+\overrightarrow{AC}$ có tọa độ là
	\choice
	{$(0;-9;9)$}
	{$(0;6;9)$}
	{$(0;9;-9)$}
	{\True $(0;6;-9)$}
	\loigiai{
		Ta có $\overrightarrow{AG}=(0;2;-3)$.\\
		Gọi $M$ là trung điểm của $BC$. Khi đó $\overrightarrow{AB}+\overrightarrow{AC}=2\overrightarrow{AM}=2\cdot\dfrac{3}{2}\overrightarrow{AG}=3\overrightarrow{AG}=(0;6;-9)$.
	}
\end{ex}

\begin{ex}%[Nguyễn Khánh Trọng]%[2H2H2-4]
	Trong không gian $Oxyz$, cho các điểm $A(2;-1;6)$, $B(-3;-1;-4)$, $(5;-1;0)$. Khẳng định nào sau đây đúng?
	\choice
	{$\triangle ABC$ cân}
	{$\triangle ABC$ có $3$ góc nhọn}
	{\True $\triangle ABC$ vuông}
	{$\triangle ABC$ đều}
	\loigiai{
		Ta có $\heva{&\overrightarrow{AB}=(-5;0;-10)\\&\overrightarrow{AC}=(3;0;-6)\\&\overrightarrow{BC}=(8;0;4)}\Rightarrow \heva{& AB=\sqrt{25+0+100}=5\sqrt{5}\\& AC=\sqrt{9+0+36}=3\sqrt{5}\\& BC=\sqrt{64+0+16}=4\sqrt{5}.}$\\
		Suy ra $AB^2=AC^2+BC^2$. \\
		Vậy $\triangle ABC$ vuông tại $C$.
	}
\end{ex}

\begin{ex}%[Nguyễn Khánh Trọng]%[2H2H2-4]
	Trong không gian $Oxyz$, cho hai véc-tơ $\overrightarrow{a}=\left( 2;m+1;-1 \right)$ và $\overrightarrow{b}=\left( 1;-3;2 \right)$. Với giá trị nào của $ m$ sau đây thì $\left| \overrightarrow{a}\cdot \overrightarrow{b} \right|=3$?
	\choice
	{$4$}
	{\True $0$}
	{$2$}
	{$-3$}
	\loigiai{
		Ta có 
			\allowdisplaybreaks			
			\begin{eqnarray*}
				\left| \overrightarrow{a}\cdot \overrightarrow{b} \right|=3
				&\Leftrightarrow& \left| 2\cdot 1-3\cdot \left( m+1 \right)+2\cdot (-1) \right|=3\Leftrightarrow \left| -3m-3 \right|=3\\
				&\Leftrightarrow& \left[ \begin{aligned}
					& m+1=1 \\
					& m+1=-1 \\
				\end{aligned} \right.\Leftrightarrow \left[ \begin{aligned}
					& m=0 \\
					& m=-2. \\
				\end{aligned} \right.
			\end{eqnarray*}
		}
\end{ex}

\begin{ex}%[Nguyễn Khánh Trọng]%[2H2H2-4]
	Trong không gian $Oxyz$, cho ba điểm $A(-1;-2;3)$, $B(0;3;1)$, $C(4;2;2)$. Cosin của góc $\widehat{BAC}$ bằng
	\choice
	{\True $\dfrac{9}{2\sqrt{35}}$}
	{$-\dfrac{9}{\sqrt{35}}$}
	{$-\dfrac{9}{2\sqrt{35}}$}
	{$\dfrac{9}{\sqrt{35}}$}
	\loigiai{Ta có $\overrightarrow{AB}=(1;5;-2)$, $\overrightarrow{AC}=(5;4;-1)$.\\
		Khi đó $\cos \widehat{BAC}=\left|\cos (\overrightarrow{AB},\overrightarrow{AC})\right|=\dfrac{\left|\overrightarrow{AB}\cdot \overrightarrow{AC}\right|}{\left|\overrightarrow{AB}\right|\cdot\left|\overrightarrow{AC}\right|}=\dfrac{9}{2\sqrt{35}}$.
	}	
\end{ex}



\begin{ex}%[Nguyễn Khánh Trọng]%[2H2V2-4]
	Trong không gian $Oxyz$, cho ba điểm $A(1;2;-1)$, $B(2;-1;3)$, $C(-4;7;5)$. Tọa độ chân đường phân giác trong góc $B$ của tam giác $ABC$ là
	\choice
	{$(-2;11;1)$}
	{\True $\left(-\dfrac{2}{3};\dfrac{11}{3};1\right)$}
	{$\left(\dfrac{2}{3};\dfrac{11}{3};\dfrac{1}{3}\right)$}
	{$\left(\dfrac{11}{3};-2;1\right)$}
	\loigiai{
		Ta có $AB=\sqrt{26}$, $BC=\sqrt{104}$ .\\
		Gọi $D$ là chân đường phân giác trong góc $B$ suy ra $\dfrac{DA}{DC}=\dfrac{BA}{BC}=\dfrac{1}{2}$.\\
		Mặt khác, $D$ thuộc đoạn $AC\Rightarrow \overrightarrow{AD}=\dfrac{1}{3}\overrightarrow{AC}$.\\
		Giả sử $D(x;y;z)\Rightarrow \overrightarrow{AD}=(x-1;y-2;z+1)$, $\overrightarrow{AC}=(-5;5;6)$ suy ra $$\left\{\begin{aligned}
			&x-1=-\dfrac{5}{3}\\
			&y-2=\dfrac{5}{3}\\
			&z+1=2\\
		\end{aligned}\right. \Rightarrow \left\{\begin{aligned}
			&x=-\dfrac{2}{3}\\
			&y=\dfrac{11}{3}\\
			&z=1\\
		\end{aligned}\right. \Rightarrow D\left(-\dfrac{2}{3};\dfrac{11}{3};1\right).$$}
\end{ex}

\begin{ex}%[Nguyễn Khánh Trọng]%[2H2V2-4]
	Trong không gian $Oxyz$, cho ba điểm $A(0;0;-1)$, $B(-1;1;0)$, $C(1;0;1)$. Tìm điểm $M$ sao cho $3MA^2+2MB^2-MC^2$ đạt giá trị nhỏ nhất.
	\choice
	{$M\left( \dfrac{3}{4};\dfrac{1}{2}-1 \right)$}
	{\True $M\left( -\dfrac{3}{4};\dfrac{1}{2};-1 \right)$}
	{$M\left( -\dfrac{3}{4};\dfrac{1}{2};2 \right)$}
	{$M\left( -\dfrac{3}{4};\dfrac{3}{2};-1 \right)$}
	\loigiai{
	Gọi $I$ là điểm thỏa mãn $3\overrightarrow{IA}+2\overrightarrow{IB}-\overrightarrow{IC}=\overrightarrow{0}$, suy ra
	$$\heva{
		& 3(0-x_I)+2(-1-x_I)-(1-x_I)=0 \\ 
		& 3(0-y_I)+2(1-y_I)-(0-y_I)=0 \\ 
		& 3(-1-z_I)+2(0-z_I)-(1-z_I)=0 \\ 
	} \Leftrightarrow \heva{
		& x_I=-\dfrac{3}{4} \\ 
		& y_I=\dfrac{1}{2} \\ 
		& z_I=-1}
		\Rightarrow I\left( -\dfrac{3}{4};\dfrac{1}{2};-1 \right).$$
	Ta có 
	\allowdisplaybreaks
	\begin{eqnarray*}
		3MA^2+2MB^2-MC^2&=&3\overrightarrow{MA}^2+2\overrightarrow{MB}^2-\overrightarrow{MC}^2\\
		&=&3\left( \overrightarrow{MI}+\overrightarrow{IA} \right)^2+2\left( \overrightarrow{MI}+\overrightarrow{IB} \right)^2-\left( \overrightarrow{MI}+\overrightarrow{IC} \right)^2\\
		&=&4MI^2+3IA^2+2IB^2-IC^2+2\overrightarrow{MI}\left( 3\overrightarrow{IA}+2\overrightarrow{IB}-\overrightarrow{IC} \right)\\
		&=&4MI^2+3IA^2+2IB^2-IC^2.
	\end{eqnarray*}
	Do đó yêu cầu đề bài $\Leftrightarrow MI$ nhỏ nhất $\Leftrightarrow M\equiv I$.\\
	Vậy $M\left( -\dfrac{3}{4};\dfrac{1}{2};-1 \right)$.
	}
\end{ex}

\Closesolutionfile{ans}
\indapan{6}{ans/ans-2-B8}
\cauds
\Opensolutionfile{ans}[ans/ans-2-B8-DS]
\begin{ex}%[Nguyễn Khánh Trọng]%[2H2V2-4]
	Trong không gian $Oxyz$, cho véc-tơ  $\overrightarrow{OA}=(2;-1;5)$ và điểm $B(5;-5;7)$.
	\choiceTF
	{\True Tọa độ của điểm $A$ là $(2;-1;5)$}
	{\True Gọi $C(a;b;c)$ thỏa mãn $\triangle ABC$ nhận $G(1;1;1)$ làm trọng tâm. Khi đó $a+b+c=-4$}
	{\True Nếu $A$, $B$, $M(x;y;1)$ thẳng hàng thì tổng $x+y=3$}
	{Cho $N\in(Oxy)$ để $\triangle ABN$ vuông cân tại $A$. Tổng hoành độ và tung độ của điểm $N$ bằng $3$}
	\loigiai{
		\begin{itemchoice}
			\itemch Đúng. Ta có tọa độ của điểm $A$ là $(2;-1;5)$.
			\itemch Đúng. $G$ là trọng tâm $\triangle ABC\Leftrightarrow \heva{&1=\dfrac{2+5+x_C}{3}\\&1=\dfrac{-1-5+y_C}{3}\\&1=\dfrac{5+7+z_C}{3}}\Leftrightarrow \heva{&x_C=-4\\&y_C=9\\&z_C=-9}\Rightarrow C(-4;9;-9)$.\\
			Suy ra $a+b+c=-4$.
			\itemch Đúng. Ta có $ \vec{AB}=(3;-4;2)  $, $ \vec{AM}=(x-2;y+1;-4) $.\\
			Ba điểm $ A $, $ B $, $ M $ thẳng hàng khi và chỉ khi 
					$$\vec{AM} = k\vec{AB} 
					\Leftrightarrow \heva{&x-2=k \cdot 3\\&y+1=k(-4)\\&-4=k \cdot 2}
					\Leftrightarrow \heva{&x=-4\\&y=7\\&k=-2.}$$
				Suy ra $ x+y=3 $.
			\itemch Sai. Ta có $N\in(Oxy)\Rightarrow N(x;y;0)$.\\
				Suy ra $\overrightarrow{AN}=(x-2;y+1;-5)$, $\overrightarrow{AB}=(3;-4;2)$.\\
				Ta có $\triangle ABN$ vuông cân tại $A\Leftrightarrow \heva{&AN\perp AB&(1)\\&AN= AB.&(2)}$\\
				$(1)\Leftrightarrow \overrightarrow{AN}\perp\overrightarrow{AB}\Leftrightarrow 3(x-2)-4(y+1)-10=0\Leftrightarrow 3x-4y=20\Leftrightarrow y=\dfrac{3}{4}x-5.$
				\allowdisplaybreaks
				\begin{eqnarray*}
					(2)&\Leftrightarrow& AN^2=AB^2
					\Leftrightarrow (x-2)^2+(y+1)^2+25=9+16+4\\
					&\Leftrightarrow& (x-2)^2+\left(\dfrac{3x}{4}-4\right)^2=4
					\Leftrightarrow x^2-4x+4+\dfrac{9x^2}{16}-6x+16=4\Leftrightarrow x=\dfrac{16}{5}.
				\end{eqnarray*}
				Suy ra $y=-\dfrac{13}{5}\Rightarrow N\left(\dfrac{16}{5};-\dfrac{13}{5};0\right)$. Vậy $x_N+y_N=\dfrac{3}{5}$.
		\end{itemchoice}
	}
\end{ex}

\begin{ex}%[Nguyễn Khánh Trọng]%[2H2V2-4]
	Trong không gian $Oxyz$, cho hai điểm $A(1;3;4)$ và $B(2;3;-1)$. 
	\choiceTF
	{\True Trọng tâm của tam giác $OAB$ là $G(1;2;1)$}
	{Khoảng cách giữa hai điểm $A$ và $B$ là $\sqrt{29}$}
	{\True Biết $C(a;b;c)$ là đỉnh thứ tư của hình bình hành $OABC$. Khi đó giá trị $a+b+c=-4$}
	{\True Gọi $S$ là tập hợp các giá trị của tham số $m$ để $D(m;m+2;1)$ thỏa mãn góc giữa hai véc-tơ $\overrightarrow{OA}$ và $\overrightarrow{BD}$ là $45^\circ$. Khi đó tổng các phần tử của $S$ là $\dfrac{51}{5}$}
	\loigiai{
		\begin{itemchoice}
			\itemch Đúng. Ta có $G$ là trọng tâm $\triangle OAB\Leftrightarrow\heva{&x_G=\dfrac{x_O+{{x}_A}+{{x}_B}}{3}=\dfrac{0+1+2}{3}=1\\&y_G=\dfrac{y_O+{{y}_A}+{{y}_B}}{3}=\dfrac{0+3+3}{3}=2\\&z_G=\dfrac{z_O+{{z}_A}+{{z}_B}}{3}=\dfrac{0+4-1}{3}=1.}$\\Vậy $G\left( 1;2;1 \right)$.
			\itemch Sai. Ta có $\overrightarrow{AB}=(1;0;-5)$. Khi đó $AB=\sqrt{1^2+0^2+(-5)^2}=\sqrt{26}$.
			\itemch Đúng. Ta có $\overrightarrow{OA}=(1;3;4)$ và $\overrightarrow{CB}=(2-a;3-b;-1-c)$.\\
			Vì $OABC$ là hình bình hành nên $\overrightarrow{OA}=\overrightarrow{CB}\Leftrightarrow \heva{&1=2-a\\&3=3-b\\&4=-1-c}\Leftrightarrow \heva{&a=1\\&b=0\\&c=-5.}$\\
			Vậy $a+b+c=-4$.
			\itemch Đúng. Ta có $\overrightarrow{OA}=(1;3;4)$ và $\overrightarrow{BD}=(m-2;m-1;2)$.
			\allowdisplaybreaks
			\begin{eqnarray*}
				\cos\left(\overrightarrow{OA},\overrightarrow{BD}\right)=\dfrac{\overrightarrow{OA}\cdot\overrightarrow{BD}}{\left|\overrightarrow{OA}\right|\cdot\left|\overrightarrow{BD}\right|}
				&\Leftrightarrow& \dfrac{\sqrt{2}}{2}=\dfrac{1\cdot(m-2)+3\cdot(m-1)+4\cdot2}{\sqrt{1+9+16}\cdot\sqrt{(m-2)^2+(m-1)^2+4}}\\
				&\Leftrightarrow& \sqrt{13}\cdot\sqrt{2m^2-6m+9}=4m+3\\
				&\Leftrightarrow& \heva{&m\ge-\dfrac{3}{4}\\&26m^2-78m+117=16m^2+24m+9}\\
				&\Leftrightarrow& \heva{&m\ge-\dfrac{3}{4}\\&10m^2-102m+108=0}\\
				&\Leftrightarrow& \heva{&m\ge-\dfrac{3}{4}\\&\hoac{&m=9\\&m=\dfrac{6}{5}}}\Leftrightarrow \hoac{&m=9\\&m=\dfrac{6}{5}.}
			\end{eqnarray*}
			Suy ra $S=\left\{\dfrac{6}{5};9\right\}$. Do đó tổng các phần tử của $S$ là $\dfrac{6}{5}+9=\dfrac{51}{5}$.
		\end{itemchoice}
	}
\end{ex}

\begin{ex}%[Nguyễn Khánh Trọng]%[2H2V2-4]
	Trong không gian $Oxyz$, cho $\triangle ABC$ có $A(1;2;3)$, $B(-2;4;4)$ và $C(4;0;5)$.
	\choiceTF
	{\True Tọa độ véc-tơ $\overrightarrow{AB}=(-3;2;1)$}
	{Trung điểm $I$ của đoạn thẳng $BC$ có tọa độ là $(1;2;4)$}
	{\True Điểm $N$ thỏa mãn $\overrightarrow{AN}=3\overrightarrow{BN}+\overrightarrow{NC}$ có tọa độ là $(-11;10;4)$}
	{\True Cho $G$ là trọng tâm $\triangle ABC$. Biết điểm $M$ nằm trên mặt phẳng $(Oxy)$ sao cho độ dài đoạn thẳng $GM$ ngắn nhất. Khi đó độ dài đoạn thẳng $GM=4$}
	\loigiai{
		\begin{itemchoice}
			\itemch Đúng. Ta có $\overrightarrow{AB}=(-3;2;1)$.
			\itemch Sai. Ta có $I$ là trung điểm của $BC\Rightarrow  I\left(\dfrac{-2+4}{2};\dfrac{4+0}{2};\dfrac{4+5}{2}\right)=\left(1;2;\dfrac{9}{2}\right)$.
			\itemch Đúng. Gọi $N(x;y;z)$. \\
			Ta có $\overrightarrow{AN}=3\overrightarrow{BN}+\overrightarrow{NC}\Leftrightarrow \heva{&x-1=3(x+2)+4-x\\&y-2=3(y-4)+0-y\\&z-3=3(z-4)+5-z}\Leftrightarrow\heva{&x=-11\\&y=10\\&z=4.}$\\
			Vậy $N(-11;10;4)$.
			\itemch Đúng. $G$ là trọng tâm tam giác $ABC$ nên $G(1;2;4)$.\\	
			Độ dài $GM$ ngắn nhất $\Leftrightarrow M$ là hình chiếu của $G$ trên $(Oxy)\Leftrightarrow M(1;2;0)$.\\
			Vậy $GM=\sqrt{0^2+0^2+(-4)^2}=4$.
		\end{itemchoice}
		
	}
\end{ex}

\begin{ex}%[Nguyễn Khánh Trọng]%[2H2V2-4]
	Trong không gian $Oxyz$, cho hai điểm $M(2;3;-1)$, $N(-1;1;1)$.
	\choiceTF
	{Hình chiếu của điểm $M$ trên trục $Oy$ có tọa độ là $(-2;3;1)$}
	{\True Gọi $E$ là điểm đối xứng của điểm $M$ qua $N$. Tọa độ của điểm $E$ là $(-4;-1;3)$}
	{\True Cho $P(1;m-1;3)$. Tam giác $MNP$ vuông tại $N$ khi và chỉ khi $m=1$}
	{Điểm $I(a;b;c)$ nằm trên mặt phẳng $(Oxy)$ thỏa mãn $T=\left|3\overrightarrow{IM}-\overrightarrow{IN}\right|$ đạt giá trị nhỏ nhất. Khi đó $2a+b+c=9$}
	\loigiai{ 
		\begin{itemchoice}
			\itemch Sai. Hình chiếu của điểm $M$ trên trục $Oy$ có tọa độ là $(0;3;0)$.
			\itemch Đúng. Vì $N$ là trung điểm của $ME\Leftrightarrow \heva{&-1=\dfrac{2+x_E}{2}\\&1=\dfrac{3+y_E}{2}\\&1=\dfrac{-1+z_E}{2}}\Leftrightarrow \heva{&x_E=-4\\&y_E=-1\\&z_E=3}\Rightarrow E(-4;-1;3)$.
			\itemch Đúng. Ta có $\overrightarrow{NM}=(3;2;-2)$; $\overrightarrow{NP}=(2;m-2;2)$.\\
			$\triangle MNP$ vuông tại $N\Leftrightarrow\overrightarrow{NM}\cdot\overrightarrow{NP}=0\Leftrightarrow 3\cdot 2+2\cdot(m-2)+(-2)\cdot 2=0 \Leftrightarrow m = 1$.
			\itemch Sai. Gọi $J(x;y;z)$ thỏa $3\overrightarrow{JM}-\overrightarrow{JN}=\overrightarrow{0}\Leftrightarrow \heva{&3(2-x)-(-1-x)=0\\&3(3-y)-(1-y)=0\\&3(-1-z)-(1-z)=0}\Leftrightarrow \heva{&x=\dfrac{7}{2}\\&y=4\\&z=-2.}$\\
			Suy ra $J\left(\dfrac{7}{2};4;-2\right)$.\\
			Khi đó $T=\left|3\overrightarrow{IM}-\overrightarrow{IN}\right|=\left|3\overrightarrow{IJ}+3\overrightarrow{JM}-\overrightarrow{IJ}-\overrightarrow{JN}\right|=\left|2\overrightarrow{IJ}\right|=2IJ$.\\
			$T$ đạt giá trị nhỏ nhất khi và chỉ khi $I$ là hình chiếu của $J$ trên $(Oxy)\Leftrightarrow I\left(\dfrac{7}{2};4;0\right)$.\\
			Vậy $a=\dfrac{7}{2}$; $b=4$; $c=0$. Suy ra $2a+b+c=11$.
		\end{itemchoice}
		
		
	}
\end{ex}

\Closesolutionfile{ans}
\indapan{3}{ans/ans-2-B8-DS}

\Opensolutionfile{ans}[ans/ans-2-B8-KQ]
\caukq

\begin{ex}%[Nguyễn Khánh Trọng]%[2H2H2-2]
	Trong không gian $Oxyz$, cho hình hộp $MNPQ.M'N'P'Q'$ với $M(1;0;0)$, $N(2;-1;1)$, $Q(0;1;0)$, $M'(1;2;1)$. Biết tọa độ điểm $P'(a;b;c)$. Khi đó giá trị tổng $a+b+c$ bằng bao nhiêu?
	\shortans{$5$}
	\loigiai{
		Ta có $\overrightarrow{MN}=(1;-1;1)$; $\overrightarrow{MQ}=(-1;1;0)$; $\overrightarrow{MM'}=(0;2;1)$.\\ 
		Do đó $\overrightarrow{MP'}=\overrightarrow{MN}+\overrightarrow{MQ}+\overrightarrow{MM'}=(0;2;2)$.\\
		Suy ra $\heva{&a-1=0\\&b-0=2\\&c-0=2}\Rightarrow a=1$; $b=2$; $c=2$.\\ 
		Vậy $a+b+c=5$.
	}
\end{ex}

\begin{ex}%[Nguyễn Khánh Trọng]%[2H2H2-4]
	Trong không gian $Oxyz$, cho 3 điểm $M\left( 2;3;-1 \right)$, $N\left( -1;1;1 \right)$, $P\left( 1;m-1;2 \right)$. Với giá trị nào của $m$ thì tam giác $MNP$ vuông tại $N$?
	\shortans{$0$}
	\loigiai{
	Ta có $\overrightarrow{MN}=\left( -3;-2;2 \right)$, $\overrightarrow{NP}=\left( 2;m-2;1 \right)$.\\
	Tam giác $MNP$ vuông tại $N$ khi và chỉ khi $$\overrightarrow{MN}\cdot\overrightarrow{NP}=0\Leftrightarrow -6-2\left( m-2 \right)+2\cdot 1=0\Leftrightarrow m=0.$$
	}
\end{ex}

\begin{ex}%[Nguyễn Khánh Trọng]%[2H2V2-4]
	Trong không gian $Oxyz$, cho véc-tơ $\vec{u}=(1;1;-2)$, $\vec{v}=(1;0;m)$. Với $m=m_0$ thì góc giữa hai $\vec{u}$, $\vec{v}$ bằng $45^{\circ}$. Khi đó giá trị $10m_0$ bằng bao nhiêu (kết quả làm tròn đến hàng phần chục)?
	\shortans{$-4{,}5$}
	\loigiai{
		\allowdisplaybreaks
		\begin{eqnarray*}
			\text{Ta có } \cos(\vec{u},\vec{v})= \dfrac{\vec{u}\cdot\vec{v}}{\left|\vec{u}\right|\cdot\left|\vec{v}\right|}&\Leftrightarrow&\dfrac{1-2m}{\sqrt{1^2+1^2+(-2)^2}\cdot\sqrt{1^2+m^2}}=\dfrac{1-2m}{\sqrt{6}\cdot\sqrt{1+m^2}}=\dfrac{\sqrt 2}{2}\\
			&\Leftrightarrow& 1-2m=\sqrt{3}\cdot\sqrt{1+m^2}
			\Leftrightarrow \heva{&1-2m\ge0\\&4m^2-4m+1=3+3m^2}\\
			&\Leftrightarrow& \heva{&m\le\dfrac{1}{2}\\&\hoac{&m=2-\sqrt{6}\\&m=2+\sqrt{6}}}\Leftrightarrow m=2-\sqrt{6}.
		\end{eqnarray*}
		Vậy $10m_0=10\left(2-\sqrt{6}\right)\approx -4{,}5$.
	}
\end{ex}

\begin{ex}%[Nguyễn Khánh Trọng]%[2H2V2-4]
	Trong không gian $Oxyz$, cho ba điểm $A(3;1;0)$, $B(2;1;-1)$, $C(x;y;-1)$ sao cho tam giác $ABC$ là tam giác vuông cân tại $A$. Khi đó tổng $x+y$ bằng bao nhiêu?
	\shortans{$5$}
	\loigiai{
	Ta có $\overrightarrow{AB}=\left(-1;0;-1 \right)$, $\overrightarrow{AC}=\left( x-3;y-1;-1 \right)$.\\
	Tam giác $ABC$ là tam giác vuông cân tại $A$ khi và chỉ khi 
	$$\heva{
		& \overrightarrow{AB}\cdot\overrightarrow{AC}=0 \\ 
		& \left| \overrightarrow{AB} \right|=\left| \overrightarrow{AC} \right|}
	\Leftrightarrow \heva{
		& -x+4=0 \\ 
		& (x-3)^2+(y-1)^2+1=2}
	\Leftrightarrow \heva{
		& x=4 \\ 
		& (y-1)^2=0}
	\Leftrightarrow \heva{
		& x=4 \\ 
		& y=1.}$$
		Vậy $x+y=5$.
	}
\end{ex}

\begin{ex}%[Nguyễn Khánh Trọng]%[2H2C2-4]
	Trong không gian $Oxyz$ cho các điểm $A(1;0;0)$, $B(3;2;4)$, $C(0;5;4)$. Xét điểm $M$ thuộc mặt phẳng $(Oxy)$ sao cho $P=MA^2+MB^2+2MC^2$ đạt giá trị nhỏ nhất. Tìm giá trị nhỏ nhất $P_{\min}$ đó.
	\shortans{$66$}
	\loigiai{
		Lấy điểm $I$ sao cho $\vec{IA}+\vec{IB}+2\vec{IC}=\vec{0}\Rightarrow \heva{
			& x_I=\dfrac{x_A+x_B+2x_C}{4}=1 \\
			& y_I=\dfrac{y_A+y_B+2y_C}{4}=3 \\
			& z_I=\dfrac{z_A+z_B+2z_C}{4}=3}\Rightarrow I(1;3;3)$.\\
		Ta có
		\allowdisplaybreaks
		\begin{eqnarray*}
			P&=&\left(\vec{MA}\right)^2+\left(\vec{MB}\right)^2+2\left(\vec{MC}\right)^2=\left( \vec{MI}+\vec{IA} \right)^2+\left( \vec{MI}+\vec{IB} \right)^2+2\left( \vec{MI}+\vec{IC} \right)^2\\
			&=&4MI^2+IA^2+IB^2+IC^2+2\vec{MI}\left( \vec{IA}+\vec{IB}+2\vec{IC} \right)\\
			&=&4MI^2+IA^2+IB^2+IC^2
		\end{eqnarray*}
		$P$ đạt giá trị nhỏ nhất $\Leftrightarrow M$ là hình chiếu của $I$ lên mặt phẳng $(Oxy)\Leftrightarrow M(1;3;0)$.\\
		Khi đó $\overrightarrow{MI}=(0;0;3)$, $\overrightarrow{IA}=(0;-3;-3)$, $\overrightarrow{IB}=(2;-1;1)$, $\overrightarrow{IC}=(-1;2;1)$.\\ 
		Suy ra $MI=3$, $IA=3\sqrt{2}$, $IB=\sqrt{6}$, $IC=\sqrt{6}$.\\
		Vậy $P_{\min}=4\cdot3^2+\left(3\sqrt{2}\right)^2+\left(\sqrt{6}\right)^2+\left(\sqrt{6}\right)^2=66$.
		}
\end{ex}

\begin{ex}%[Nguyễn Khánh Trọng]%[2H2C2-4]
	Trong không gian $Oxyz$, cho véc-tơ $\overrightarrow{a}=(1;-1;0)$ và hai điểm $A(-4;7;3)$, $B(4;4;5)$. Hai điểm $M$, $N$ thay đổi thuộc mặt phẳng $(Oxy)$ sao cho $\overrightarrow{MN}$ cùng hướng với $\overrightarrow{a}$ và $MN=5\sqrt{2}$. Giá trị lớn nhất của $\left| AM-BN \right|$ bằng bao nhiêu (kết quả làm tròn đến hàng phần trăm)? 
	\shortans{$4{,}12$}
	\loigiai{
	\immini{
			Vì $\overrightarrow{MN}$ cùng hướng với $\overrightarrow{a}$ nên tồn tại số thực $k>0$ sao cho $$\overrightarrow{MN}=k\overrightarrow{a}\Leftrightarrow MN=\left| k \right|\cdot\left| \overrightarrow{a} \right|\Leftrightarrow \left| k \right|=5\Leftrightarrow k=5.$$ 
			Gọi $K\left( x;y;z \right)$ thỏa mãn
			$\overrightarrow{AK}=\overrightarrow{MN}$, ta có 
			$$\heva{& x+4=5 \\ & y-7=-5 \\ & z-3=0}
			\Leftrightarrow \heva{& x=1 \\ & y=2 \\ & z=3}
			\Rightarrow K\left( 1;2;3 \right).$$
		}{\begin{tikzpicture}[line join=round,line cap=round,>=stealth,scale=0.6,font=\footnotesize]
				\tikzset{every node/.style={outer sep=-0.5mm}}
				\foreach \x/\y/\n in {-0.75/0/a,8/0/d,-2/-2/b,-1/2/A,0.25/-1/M,6/1/B} \coordinate (\n) at (\x,\y);
				\coordinate (c) at ($(b)+(d)-(a)$); 
				\coordinate (x) at ($(a)+(6,1.75)$);
				\coordinate (y) at ($(x)+(1.5,0)$); 
				\coordinate (z) at ($(x)!3!(y)$);
				\coordinate (N) at ($(M)+(z)-(x)$); 
				\coordinate (K) at ($(A)+(N)-(M)$); 
				\draw (a)--(b)--(c)--(d)--(a);
				\draw[thick,arrow style] (A)--(K);
				\draw[thick,arrow style] (M)--(N);
				\draw[thick,blue] (A)--(M) (K)--(N);
				\draw[thick,red] (B)--(N);
				\draw[thick,arrow style] (x)--(y) node[midway,sloped,above]{$\vec{a}$};
				\foreach \n/\l in {K/above,N/below,M/below,A/above,B/above} \fill (\n) node[\l]{$\n$} circle (1.5pt);
				\draw pic[thin,draw,"$Oxy$",angle radius=1.2cm] {angle= c--b--a};
			\end{tikzpicture}}
		\noindent Suy ra $K$ và $B$ nằm cùng phía đối với $\left( Oxy \right)$ và $AM=KN$ (do $AKNM$ là hình bình hành). Khi đó
		$$\left| AM-BN \right|=\left| KN-BN \right|\le KB=\sqrt{17}.$$
		Dấu \lq\lq$=$\rq\rq\ xảy ra $\Leftrightarrow K$, $N$, $B$ thẳng hàng.\\
		Vậy giá trị lớn nhất của $\left| AM-BN \right|$ bằng $\sqrt{17}\approx 4{,}12$.
	}
\end{ex}

\Closesolutionfile{ans}
\indapan{6}{ans/ans-2-B8-KQ}